\documentclass[10pt,a4paper]{amsart}
\usepackage[margin=19mm]{geometry}
\usepackage{amsmath,amssymb,mathtools}
\usepackage[scr=rsfs,cal=euler]{mathalfa}
\usepackage{xfrac}
\usepackage[unicode,hidelinks,bookmarks=false]{hyperref}
\newcommand{\ie}{i.e.\ }
\newcommand{\cf}{cf.\ }
\newcommand{\resp}{respectively\ }
\newcommand{\Subsection}{\subsection}
\providecommand{\nobreakdash}{\nobreak\hbox{-}}
\setlength{\emergencystretch}{2em}
\multlinegap0pt
\usepackage{amssymb}
\usepackage{mathtools}
\usepackage{float}
\usepackage[shortlabels]{enumitem}
\usepackage{xcolor}
\newtheorem{lemma}{Lemma}
\newtheorem{proposition}{Proposition}
\usepackage{cleveref}
\crefname{lemma}{Lemma}{Lemmas}
\crefname{proposition}{Proposition}{Propositions}
\newcommand{\Int}{\operatorname{Int}}
\title{Weak Polynomial Completeness Does Not Imply \\ Almost Polynomial Completeness}
\author{Viet-Hoang Tran, Tan N. Nguyen}
\date{Source: arXiv:2609.05891v1 (CC BY 4.0)}
\begin{document}
\maketitle

\noindent\emph{Source and licence.} Weak Polynomial Completeness Does Not Imply \\ Almost Polynomial Completeness. Version arXiv:2609.05891v1 (CC BY 4.0), distributed by arXiv under the Creative Commons Attribution 4.0 International licence, http://creativecommons.org/licenses/by/4.0/, as recorded in the arXiv licence metadata for this exact version and shown on the abstract page https://arxiv.org/abs/2609.05891v1. Full licence notice: \texttt{licenses/arxiv-2609.05891-CC-BY-4.0.txt}.\par\smallskip

\noindent\textit{Editorial scope.} Counted unit: the article's proposition prop:invariant (source0.tex lines 465--475). The article's complete original proof of it is delivered with it (source0.tex lines 477--547, one \texttt{proof} environment of 71 lines). No mathematics is written, repaired or re-derived: the statement and the proof are the article's own lines, in the article's own notation, and no retained line is substituted or re-flowed.  Retained uncounted as the source-local context the proof needs: lines 331--336; lines 391--395; lines 455--461.  Nothing was omitted from any retained span, so the package carries no omission record.  No retained line contains a citation, so the wrapper carries no bibliography and no bibliography entry.  No retained line contains a figure environment, an image-inclusion command or drawing code, so no image is delivered and the package declares no asset dependency.  Editorial formatting only, in addition: an AMS \texttt{amsart} preamble that restates the article's own declarations and macros used by the retained lines, namely the environments \texttt{lemma}, \texttt{proposition} on separate counters, together with the standard packages those lines need.  The retained environments print in this document as Lemma 1, Proposition 1 in order of appearance, whereas the article prints the same items with the numbering of its own full document.  The complete native source of the article as distributed in the arXiv e-print for this exact version is bundled unchanged as \texttt{sources/209472/source0.tex}.
\begin{lemma}[Unary derivative lemma]\label{lem:derivative}
For every \(f\in\Int(D)\) and every \(a\in D\), one has
\[
 f'(a)\in V.
\]
\end{lemma}

\begin{equation}
 A_{r+1}
  =D\big[A_r,\ f(G):f\in\Int(D),\ G\in A_r\big]\subseteq C.
 \label{eq:Ar}
\end{equation}

\begin{equation}
 E=\left\{
 G\in C:
 G_U(0,0),\ G_W(0,0),\ G_{UW}(0,0)\in V
 \right\}.
 \label{eq:E}
\end{equation}

\begin{proposition}\label{prop:invariant}
The set \(E\) is a \(D\)-subalgebra of \(C\), it contains
\(D[U,W]\), and it is closed under every unary operation
\[
 G\longmapsto f(G),\qquad f\in\Int(D).
\]
Consequently,
\[
 A\subseteq E.
\]
\end{proposition}

\begin{proof}
Closure under addition and multiplication by elements of \(D\) follows
directly from linearity of formal differentiation and \(D\subseteq V\).
Let \(G,H\in E\).  The product rules give
\[
 (GH)_U=G_UH+GH_U,\qquad
 (GH)_W=G_WH+GH_W
\]
and
\begin{equation}
 (GH)_{UW}
  =G_{UW}H+G_UH_W+G_WH_U+GH_{UW}.
 \label{eq:product-mixed}
\end{equation}
At the origin, \(G(0,0)\) and \(H(0,0)\) lie in \(D\subseteq V\),
while all derivatives occurring on the right sides lie in \(V\) by
the definition of \(E\).  Every displayed derivative of \(GH\) at the
origin therefore belongs to \(V\).  This proves closure under
multiplication.  Closure under additive inverses also follows from
linearity of differentiation.

For a polynomial
\(G=\sum_{i,j}a_{ij}U^iW^j\in D[U,W]\), the three derivatives at the
origin are \(a_{10}\), \(a_{01}\), and \(a_{11}\), respectively (with
a missing coefficient interpreted as \(0\)).  In particular, they
belong to \(D\subseteq V\).
Thus
\[
 D[U,W]\subseteq E.
\]
In particular, \(E\) contains \(D\); together with the preceding
closure properties, this proves that \(E\) is a \(D\)-subalgebra of
\(C\).

It remains to prove unary closure.  Take \(G\in E\) and
\(f\in\Int(D)\), and put \(a=G(0,0)\in D\).  The formal chain rule
gives
\begin{equation}
 \begin{aligned}
 (f(G))_U&=f'(G)G_U,\\
 (f(G))_W&=f'(G)G_W,\\
 (f(G))_{UW}
   &=f''(G)G_UG_W+f'(G)G_{UW}.
 \end{aligned}
 \label{eq:chain}
\end{equation}
The base field has characteristic \(2\).  If
\(f(Z)=\sum_j c_jZ^j\), then
\[
 f''(Z)=\sum_{j\geq2}j(j-1)c_jZ^{j-2}=0,
\]
because every integer \(j(j-1)\) is even.  Evaluating
\eqref{eq:chain} at the origin therefore yields
\[
 \begin{aligned}
 (f(G))_U(0,0)&=f'(a)G_U(0,0),\\
 (f(G))_W(0,0)&=f'(a)G_W(0,0),\\
 (f(G))_{UW}(0,0)&=f'(a)G_{UW}(0,0).
 \end{aligned}
\]
By \cref{lem:derivative}, \(f'(a)\in V\); by hypothesis, the other
factors lie in \(V\).  All three products lie in \(V\), and so
\(f(G)\) satisfies the derivative conditions in \eqref{eq:E}.
Moreover \(G\in C\) implies \(f(G)\in C\), as observed immediately
after \eqref{eq:Ar}.  Hence \(f(G)\in E\).

Since \(A_0=D[U,W]\subseteq E\), the unary closure just proved and the
fact that \(E\) is a \(D\)-algebra imply successively that
\(A_r\subseteq E\) for every \(r\).  Taking the union gives
\(A\subseteq E\).
\end{proof}

\end{document}
